\chapter{Admission общего relative-Hodge auxiliary-edge parent}
\label{ch:v10-h15-r2-relative-hodge-auxiliary-edge-admission}

Условный relative parent требует две изоморфные восьмимерные копии:
физическую $\Sigma$ и auxiliary-копию $A_\Sigma$. Проверим одновременно
типизацию второй копии и Hodge-структуру двухузлового графа.

Унаследованный 36-мерный fixed-point сектор имеет только веса
$0,\pm4,\pm6$, тогда как требуемая копия несёт
\begin{equation}
 (3,-3,7,1,-1,-7,3,-3).
\end{equation}
Матрица совпадений весов равна нулю, а полный intertwiner constraint имеет
rank/nullity $288/0$. Поэтому существующий auxiliary-сектор не содержит ни
одного гиперзаряд-эквивариантного направления нужного типа.

Минимальное условное расширение есть
\begin{equation}
 \mathcal H_{\rm rel}=\Sigma\otimes\mathbb C^2,
 \qquad \dim_{\mathbb R}\mathcal H_{\rm rel}=16.
\end{equation}
На нём gauge generator равен $\operatorname{diag}(Q,Q)$, а self-adjoint
odd edge меняет две копии местами. Ориентированная incidence-карта
$(I_8,-I_8)$ имеет rank $8$ и порождает
\begin{equation}
 P_{\rm rel}=\frac12
 \begin{pmatrix}I_8&-I_8\\-I_8&I_8\end{pmatrix}.
\end{equation}
Этот оператор идемпотентен, имеет rank $8$ и inertia
$(0_-,8_0,8_+)$. Он коммутирует с gauge generator и Real-структурой;
odd edge точно антикоммутирует с grading. Поэтому условная архитектура
закрывает все шесть admission-слотов и сохраняет exact $Q$ readout.

Но это именно расширение. Inherited Hodge-selector имеет rank $0$, а
equivariant Hom из требуемого auxiliary-модуля в fixed-point сектор равен
нулю. Минимальный дефицит составляет восемь вещественных направлений.

\begin{theorem}[Минимальный условный doubled carrier]
$\Sigma\otimes\mathbb C^2$ является минимальным Real-, grading- и
gauge-совместимым носителем общего relative-Hodge auxiliary edge. Его
Hodge projector каноничен и положителен. Однако в текущем fixed-point
parent отсутствует вся вторая весовая копия, поэтому admission физически
равен лишь $3/6$.
\end{theorem}

\begin{nogocriterion}[Размерность без совпадения весов не является embedding]
Нельзя использовать 36-мерную fixed-point algebra как auxiliary carrier
только потому, что её размерность больше восьми. Нулевой weight-match и
нулевой equivariant Hom запрещают такое вложение; требуется независимое
происхождение новой копии $\Sigma$ и её graph-Hodge edge.
\end{nogocriterion}

\begin{protocol}\begin{enumerate}
\item fixed-point/required weight-match rank: \textbf{$0$};
\item equivariant Hom dimension: \textbf{$0$};
\item minimum extension dimension: \textbf{$8$};
\item conditional common carrier dimension: \textbf{$16$};
\item oriented incidence rank: \textbf{$8$};
\item Hodge projector rank/nullity: \textbf{$8/8$};
\item Hodge inertia: \textbf{$(0_-,8_0,8_+)$};
\item gauge/Real/grading compatibility: \textbf{точно};
\item inherited/conditional admission: \textbf{$3/6$ и $6/6$};
\item следующий гейт: \textbf{аудит происхождения auxiliary edge}.
\end{enumerate}\end{protocol}

Машинный сертификат сохранён в
\path{s2t_v10_particle_wrinkle_dislocation_callias_equal_charge_twist_m4_h15_r2_relative_hodge_auxiliary_edge_common_parent_admission_gate_results.json}.